\subsection{Rational approximation and an infinitely supported unit equation}

\begin{statement}{A.42}{Proved}{The rational unit-equation graph is closed and discrete}
Give $\widehat{\Q^\times}$ the inverse-limit topology of A.10,
and put
\[
 \mathcal R=
 \{(\kappa(t),\kappa(1-t)):t\in\Q\setminus\{0,1\}\}
       \subset\widehat{\Q^\times}^{\,2}.
\]
This subset is closed and discrete. Every Cauchy net in
$\mathcal R$ is eventually constant. For a pair
$(A,B)\notin\mathcal R$, there is a positive integer $N$
such that no rational $t\ne0,1$ has both
$[t]_N=A_N$ and $[1-t]_N=B_N$ in
$\Q^\times/(\Q^\times)^N$.
\end{statement}
\begin{proof}
Fix classes $\alpha,\beta\in\Q^\times/(\Q^\times)^4$
and rational representatives $a,b\ne0$. The smooth projective
plane curve
\[
 C_{a,b}:aX^4+bY^4=Z^4
 \tag{A.42.1}
\]
has genus three: its three partial derivatives cannot vanish
simultaneously at a projective point, and the smooth
plane-quartic genus formula gives $(4-1)(4-2)/2=3$.
On $C_{a,b}^{\circ}=C_{a,b}\cap\{XYZ\ne0\}$ the morphism
\[
 C_{a,b}^{\circ}\longrightarrow U,
 \qquad [X:Y:Z]\longmapsto a(X/Z)^4
 \tag{A.42.2}
\]
has $1-t=b(Y/Z)^4$. Conversely, a rational $t$ with
$[t]_4=\alpha$ and $[1-t]_4=\beta$ gives nonzero rational
$x,y$ satisfying $t=ax^4$, $1-t=by^4$, hence the rational
point $[x:y:1]$ of (A.42.1). Faltings's theorem
\cite[Satz 7, printed p.~365]{AgentFaltings1983} makes
$C_{a,b}(\Q)$ finite. Thus each fibre of
\[
 \rho_4:\widehat{\Q^\times}^{\,2}
       \longrightarrow
       \bigl(\Q^\times/(\Q^\times)^4\bigr)^2
\]
meets $\mathcal R$ in a finite set.

The target of $\rho_4$ is discrete, so each such fibre is
open and closed. The ambient inverse limit is Hausdorff.
For a point outside $\mathcal R$, remove the finite
intersection with $\mathcal R$ from its $\rho_4$-fibre;
this gives an open neighbourhood disjoint from $\mathcal R$.
For a point in $\mathcal R$, remove the other points of
that finite intersection. This proves closedness and
discreteness, with the actual neighbourhood maps specified.

A Cauchy net is eventually in a single $\rho_4$-fibre,
and thus eventually in a finite subset of $\mathcal R$.
For each two distinct members, some power quotient separates
them; taking a common multiple of these finitely many
moduli gives a single quotient injective on the finite set.
The Cauchy condition in that quotient makes the net
eventually constant. Finally, the complement of
$\mathcal R$ is open, and the basic open neighbourhoods
are fibres of the simultaneous power-quotient maps $\rho_N$.
This proves the final assertion.
\end{proof}

In particular, a construction from rational solutions
$t_j+(1-t_j)=1$ cannot produce a new pair by making their
two global power classes Cauchy. The construction below
uses rational coefficients of locally soluble diagonal
curves; those coefficients need not themselves sum to one
in $\Q$. Their sums tend to one in each $\Q_p$, and we
prove separately that both multiplicative sequences
converge in the topology of A.10.

% Insert as A.43.  The construction in A.44 has the valuations in A.43.6.
\begin{statement}{A.43}{Proved}{The weight-three test on the twenty-four integral lifts}
Put $E_3=E_U/\Gamma_4\Pi$, where
$\Gamma_4\Pi=\overline{[\Pi,\Gamma_3\Pi]}$.  With the generators
and tangent of A.38, set
\[
 C=[x,y],\qquad D=[C,x],\qquad F=[C,y].
\]
The kernel of $E_3\to E_2$ is
$\widehat\Z(3)D\oplus\widehat\Z(3)F$.
For a class-two section $s(g)=(a(g),b(g),h(g);g)$ as in A.38,
write $\mathbf o_3(s)=(o_D(s),o_F(s))$ for its actual lifting
obstruction in $H^2_{\rm cts}(G_\Q,\widehat\Z(3))^2$.
Set $\varepsilon=(\chi-1)/2$ and $\eta=(\chi^2-1)/24$.
Its components are represented by the continuous two-cocycles
\[
 \begin{split}
 W_D&=-(a+\varepsilon)\cup h-\binom a2\cup b,\\
 W_F&=(b+\varepsilon)\cup(ab-h)+\binom b2\cup a-\eta\cup b.
 \end{split}
 \tag{A.43.1}
\]
Here $ab$ and $\binom a2$ are pointwise one-cochains with
coefficient weight two.  In particular all the expressions are
integral, including at two.  The twenty-four classes in A.40
have the exact obstruction values
\[
 \mathbf o_3(s+k\eta)=\mathbf o_3(s)
   -k\bigl((A+\kappa(-1))\cup[\eta],
           (B+\kappa(-1))\cup[\eta]\bigr),
                    \qquad k\in\Z/24.
 \tag{A.43.2}
\]

Let $(A,B)\in\mathcal G$, with ordinary coordinates $x_q$, and
let $\delta_q$ be its comparison class from (A.40.16).  For
$q\ne\ell$ denote the Frobenius coordinate of its $\ell$-part
by
\[
 d_{q,\ell}\in\Z_\ell/(q^2-1)\Z_\ell
       =H^1(G_{\Q_q},\Z_\ell(2)).
\]
Under the standard tame-residue isomorphism
\[
 \rho_{q,\ell}:H^2(G_{\Q_q},\Z_\ell(3))
       \xrightarrow{\ \sim\ }\Z_\ell/(q^2-1)\Z_\ell,
 \tag{A.43.3}
\]
the two localized weight-three obstructions are exactly
\[
 \rho_{q,\ell}\bigl(\mathbf o_3(s)|_{G_{\Q_q}}\bigr)
   =\bigl(-v_q(x_q)d_{q,\ell},
           -v_q(1-x_q)d_{q,\ell}\bigr).
 \tag{A.43.4}
\]
Thus their common zero locus at this place is the annihilator
of the two indicated valuation integers in the displayed finite
cyclic module.

Write $\mathscr L_i(A,B)$ for the geometric conjugacy classes
of $E_i$-sections with abelian image $(A,B)$, and $s_*$ for the
distinguished class of A.40.  On the ordinary locus with
\[
 (v_2(x_2),v_2(1-x_2))=(0,1),\qquad
 (v_3(x_3),v_3(1-x_3))=(1,0),
 \tag{A.43.5}
\]
the actual projection of section sets satisfies
\[
 \operatorname{im}\bigl(\mathscr L_3(A,B)\longrightarrow
                      \mathscr L_2(A,B)\bigr)\subseteq\{s_*\}.
 \tag{A.43.6}
\]
For the rational pair $(A,B)=(\kappa(3),\kappa(-2))$, this
image is exactly the class-two point-section class at $3$.
\end{statement}
\begin{proof}
\emph{The actual group, action and cochain representative.}
Commutator collecting generates $\Gamma_3\Pi/\Gamma_4\Pi$
by $D,F$.  To check their independence explicitly, send $x,y$
to $1+X,1+Y$ in the units of the noncommutative polynomial ring
truncated after degree three, with coefficients in
$\widehat\Z$.  This is a continuous homomorphism on the
class-three quotient.  The two commutators have expansions
\[
 \begin{split}
 D&\longmapsto 1+2XYX-YX^2-X^2Y,\\
 F&\longmapsto 1+XY^2-2YXY+Y^2X.
 \end{split}
 \tag{A.43.7}
\]
The coefficients of $X^2Y$ and $XY^2$ show that
$(u,v)\mapsto D^uF^v$ has zero kernel.  This proves the stated
basis of the kernel.  In particular every element of the geometric
class-three quotient has a unique form
\[
                         y^b x^a C^h D^uF^v.
 \tag{A.43.8}
\]
Its image in the Heisenberg model (A.38.2) is $(a,b,h)$;
the order of the first two factors is part of this identification.
Continuity follows from the profinite powers and the compact
coordinate spaces.

The arithmetic action from (A.38.6) has
$f_g\equiv C^{\eta(g)}\pmod{\Gamma_3\Pi}$, with this exact
coefficient and sign.  The coefficient calculation is
\cite[\S12.5.2, equation (12.35)]{BIR_Wickelgren2012}, using
the same tangent, generators and action
$g(y)=f_g^{-1}y^{\chi(g)}f_g$.
Writing $\lambda=\chi(g)$, the action on (A.43.8) is therefore
determined by
\[
 \begin{split}
 g(x)&=x^\lambda,\qquad
 g(y)=y^\lambda F^{-\lambda\eta(g)},\\
 g(C)&=C^{\lambda^2}
       D^{-\lambda\binom\lambda2}
       F^{-\lambda\binom\lambda2},\\
 g(D)&=D^{\lambda^3},\qquad g(F)=F^{\lambda^3}.
 \end{split}
 \tag{A.43.9}
\]
These formulas also verify the asserted cyclotomic action on
the kernel.  All binomial expressions are continuous
integer-valued functions on $\widehat\Z$.

The coordinate lift of $s$ to the geometric class-three quotient
is $\widetilde s(g)=y^{b(g)}x^{a(g)}C^{h(g)}$.  Its defect
\[
 \widetilde s(g)\,g(\widetilde s(t))\,
                       \widetilde s(gt)^{-1}
\]
has coordinates $(R_D,R_F)$ on the central kernel.  Applying
(A.43.9) and collecting the group words gives the exact cochain
identities
\[
 \begin{split}
 R_D&=W_D-\mathrm d(ah),\\
 R_F&=W_F+\mathrm d\bigl(a\binom b2-bh\bigr).
 \end{split}
 \tag{A.43.10}
\]
For completeness this collecting calculation can be performed
inside the same degree-three polynomial ring: use
\[
 \begin{split}
 (1+X)^a&=1+aX+\binom a2X^2+\binom a3X^3,\\
 C&=1+XY-YX+YX^2-XYX-XY^2+YXY,
 \end{split}
\]
and (A.43.7).  Substitute
$a(gt)=a(g)+\lambda a(t)$,
$b(gt)=b(g)+\lambda b(t)$ and
$h(gt)=h(g)+\lambda^2h(t)+a(g)\lambda b(t)$.
The degree-one and degree-two terms cancel.  The negative
$X^2Y$ coefficient and the $XY^2$ coefficient are the two
expressions in (A.43.10); these determine the central defect
by (A.43.7).  This proves (A.43.10) as an identity of
integer-valued polynomials, and continuity extends it to all
profinite exponents.  The symbolic program
\texttt{class3\_obstruction\_certificate.py} verifies precisely
these two identities over the rational polynomial ring, together
with (A.43.9).  Thus the computation does not discard a
finite-coefficient or continuous-cohomology correction.

The raw defect is a two-cocycle for the kernel action.  Formula
(A.43.10) proves that (A.43.1) consists of two-cocycles with the
same obstruction classes.  It is also the explicit comparison
with \cite[Proposition 17 and Theorem 19]{BIR_Wickelgren2012}:
the source's ordered coordinates $(b,a)_c=y^a x^b C^c$ become
our $(a,b,h)$ by $b_{\rm source}=a$,
$a_{\rm source}=b$, $c_{\rm source}=h$.
The obstruction is zero exactly on the image of the actual
map of section sets induced by $E_3\to E_2$: a continuous
primitive for the defect changes the last two coordinates of
the lift (A.43.8) to a homomorphism, and a section conversely
supplies such a primitive.

Replacing $h$ by $h+c$ for a continuous central cocycle gives,
by direct subtraction in (A.43.1),
\[
 \mathbf o_3(s+c)-\mathbf o_3(s)
       =-\bigl((A+\kappa(-1))\cup[c],
               (B+\kappa(-1))\cup[c]\bigr).
 \tag{A.43.11}
\]
Here $[\varepsilon]=\kappa(-1)$ by the compatible roots of
$-1$.  A.40 now proves (A.43.2), with no further choices of
defining systems for Massey products.

\emph{The exact local residue.}
The ordinary point section at $x_q$ has zero obstruction and
differs from $s|_{G_{\Q_q}}$ by $\delta_q$ after the unique
abelian alignment of A.40.  Applying (A.43.11) locally gives
\[
 \mathbf o_3(s)|_{G_{\Q_q}}
   =\bigl((A_q+\kappa_q(-1))\cup\delta_q,
           (B_q+\kappa_q(-1))\cup\delta_q\bigr).
 \tag{A.43.12}
\]
For $q\ne\ell$ every integral weight-two cocycle is unramified
by (A.40.6).  The tame Hochschild--Serre calculation identifies
$H^2(G_{\Q_q},\Z_\ell(3))$ with
$H^1(\operatorname{Gal}(\bar\F_q/\F_q),\Z_\ell(2))$:
prime-to-$\ell$ inertia has zero higher $\ell$-cohomology,
tame $\ell$-inertia has cohomological dimension one, and the
procyclic quotient has cohomological dimension one.  The same
calculation on finite coefficients, followed by inverse limit,
is valid for continuous integral cohomology since the finite
local $H^1$ groups are finite and their $\varprojlim^1$ is zero.
Frobenius evaluation now gives the isomorphism (A.43.3).

We fix its sign and normalization by the cochain residue from
A.27, now with coefficient weight three.  Write $T$ for a tame
inertia generator and $\Phi T\Phi^{-1}=T^q$.  For a normalized
two-cochain $v$, the residue is represented by
\[
 q^{-1}\left(v(\Phi,T)+v(\Phi T,\Phi^{-1})
       -v(\Phi,\Phi^{-1})
       -\sum_{i=1}^{q-1}v(T^i,T)\right).
 \tag{A.43.13}
\]
The telescoping computation of A.27 gives
$(q^2-1)u(T)$ on $v=\mathrm du$; this is the differential
appropriate to weight three.  The discarded inertia is handled
by exactly the coefficient-compatible averaging map of A.27.

Let $w$ be a weight-one cocycle, with $w(T)=t$ and
$w(\Phi)=\alpha$, and let $c$ be an unramified weight-two
cocycle with $c(\Phi)=d$.  In (A.43.13), the cochain
$v=w\cup c$ has first and last terms zero and middle
difference
\[
                   -\bigl(\alpha+qt\bigr)d+\alpha d=-qtd.
\]
Its residue is therefore $-td$ modulo $(q^2-1)\Z_\ell$.
The inertia values of the two weight-one Kummer classes in
(A.43.12) are precisely $v_q(x_q)$ and $v_q(1-x_q)$;
$\kappa_q(-1)$ has inertia value zero.  This proves (A.43.4)
with the exact comparison map and sign.

\emph{The two fixed valuations.}
For a class in the image of $\mathscr L_3(A,B)$, both
obstructions vanish.  Under (A.43.5), (A.43.4) at
$(q,\ell)=(2,3)$ forces $d_{2,3}=0$, and at $(3,2)$ it
forces $d_{3,2}=0$.  The isomorphism (A.40.4) says that
exactly one global class-two lift has these two comparisons
zero, namely $s_*$.  This proves (A.43.6).

For $x=3$, the full rational point section exists and gives
one element in this image.  More explicitly, take its class-two
class as $s_0$.  For $s_0+k\eta$, the two local obstruction
vectors at the indicated coefficient pairs are
\[
       (0,k/8)\in(\Z/3)^2,
       \qquad(k/3,0)\in(\Z/8)^2.
 \tag{A.43.14}
\]
They vanish together exactly for $k=0$ in $\Z/24$.
This identifies the image as the point-section class itself.
\end{proof}

The reduction of the integral central adjustment map at two
is especially small:
\[
 \operatorname{im}\bigl(H^1(G_\Q,\widehat\Z(2))
                \longrightarrow H^1(G_\Q,\Z/2)\bigr)
                      =\{0,\kappa_2(2)\}.
 \tag{A.43.15}
\]
Indeed $\eta\bmod2$ is one exactly when
$\chi_2\equiv\pm3\pmod8$, and zero when
$\chi_2\equiv\pm1\pmod8$.  This is the quadratic character of
$\Q(\sqrt2)$, as also computed in
\cite[Lemma 31]{BIR_Wickelgren2012}.  Formula (A.43.15)
specifies exactly which mod-two changes of central defining
system come from the continuous integral lift fiber; the
finite-coefficient choice is required to lie in this image.

The nonrational pairs constructed in A.44 have the valuations
(A.43.5).  Their possible class-three projections are therefore
confined to the single class $s_*$.  No vanishing of
$\mathbf o_3(s_*)$ for those pairs is proved here.  Equations
(A.43.1), (A.43.4) and the actual extension $E_3\to E_2$
specify the remaining calculation, including its continuous
global cohomology class.


\begin{statement}{A.44}{Proved}{Infinitely supported ordinary-local solutions of the completed unit equation}
There are recursively specified pairwise distinct primes
$q_j,r_j>3$, for $j\ge0$, such that, with $n_j=(j+4)!$,
the completed products
\[
 A=3\prod_{j\ge0}q_j^{n_j},
 \qquad B=-2\prod_{j\ge0}r_j^{n_j}
       \quad\text{in }\widehat{\Q^\times}
 \tag{A.44.1}
\]
have the following properties. At every finite prime $p$
their localizations are the Kummer classes of ordinary
nonzero numbers $x_p,y_p\in\Z_p$ with $x_p+y_p=1$.
Thus $(A,B)\in\mathcal G$. Their valuation vectors are
\[
 \begin{array}{c|cc}
 p&v_p(A)&v_p(B)\\\hline
 3&1&0\\
 2&0&1\\
 q_j&n_j&0\\
 r_j&0&n_j\\
 \text{all other primes}&0&0.
 \end{array}
 \tag{A.44.2}
\]
Both coordinates have infinite support, so neither is the
Kummer class of a rational number. The construction uses
no local section theorem.
\end{statement}
\begin{proof}
We prove the field-intersection calculation which permits
the recursive prime choices, give an elementary finite-field
bound, and then construct the products with explicit
precision estimates.

\emph{The radical fields and their abelian subfields.}
Let $24\mid h$ and let
\[
 c=3w^8\quad\text{or}\quad c=-2w^8,
 \qquad w\in\Z_{>0},\quad\gcd(w,6)=1.
 \tag{A.44.3}
\]
Put $K=\Q(\mu_h)$, choose $\gamma^h=c$, and set
$L=K(\gamma)$. This is the splitting field of $T^h-c$
over $\Q$. We claim that
\[
                L^{\mathrm{ab}}=K,
 \tag{A.44.4}
\]
where $L^{\mathrm{ab}}$ denotes the fixed field of the
commutator subgroup of $\operatorname{Gal}(L/\Q)$.

First, $c$ is a square in $K$: the field $\Q(\mu_{24})$
contains $\sqrt3$ and $\sqrt{-2}$. It is not an $e$th
power in $K$ for any odd prime $e\mid h$. Indeed, the polynomial
$T^e-c$ is Eisenstein at $3$ in the first case and at $2$
in the second. Its real root generates a degree-$e$
subfield of $\R$ which has nonreal conjugates, so this
subfield is not Galois over $\Q$. If an $e$th root of $c$
lay in $K$, with $e\mid h$, multiplication by an element
of $\mu_e\subset K$ would put the real root in $K$.
Every subfield of the abelian extension $K/\Q$ is Galois
over $\Q$, a contradiction. Likewise $c$ is not a fourth
power in $K$: after division by $w^2$, and also by
$\zeta_8$ in the negative case, a fourth root of $c$
would put $3^{1/4}$ or $2^{1/4}$ in $K$. Each of these
positive real radicals has an Eisenstein degree-four
polynomial and nonreal conjugates.

The Kummer class of $c$ in $K^\times/(K^\times)^h$
therefore has order exactly $h/2$. To check the last
assertion explicitly, its order divides $h/2$ because
$c$ is a square. If it divided $h/(2e)$ for a prime
$e\mid h/2$, an equation $c^{h/(2e)}=z^h$ would give
$c=z^{2e}\eta$ with $\eta\in\mu_{h/(2e)}$.
Since $\mu_h\subset K$, the factor $\eta$ has a
$2e$th root in $K$. This would make $c$ an $e$th power
for odd $e$, or a fourth power for $e=2$, both excluded.
Consequently
\[
 N:=\operatorname{Gal}(L/K)=\langle\tau\rangle
      \simeq\Z/(h/2),\qquad
 \tau(\gamma)=\zeta_h^2\gamma.
\]
For $G=\operatorname{Gal}(L/\Q)$ and
$\sigma|_K(\zeta_h)=\zeta_h^u$, conjugation gives
$\sigma\tau\sigma^{-1}=\tau^u$. All
$u\in(\Z/h)^\times$ occur. The greatest common divisor
of $h/2$ and these integers $u-1$ is $2$: all are even;
one can take $u\equiv3\pmod4$, and, for every odd prime
$e\mid h$, choose a unit $u\not\equiv1\pmod e$ by
the Chinese remainder theorem. Hence
\[
                 [N,G]=N^2.
\]
Since $G/N$ is abelian, $N^2\subset[G,G]\subset N$.
The intermediate field fixed by $N^2$ is
$K(c^{1/4})$. It contains $\Q(3^{1/4},i)$ or
$\Q(2^{1/4},i)$, according to (A.44.3), because
$\zeta_8\in K$. This subfield has a nonabelian Galois
group: on the fourth roots, the rotation
$\alpha\mapsto i\alpha$ and complex conjugation do
not commute. Thus $K(c^{1/4})/\Q$ is not abelian.
As $N/N^2$ has order two, it follows that $[G,G]=N$.
This proves (A.44.4).

Here is the exact prime-selection consequence. Let $M$
be any positive multiple of $h$, and let
$u\in(\Z/M)^\times$ satisfy $u\equiv1\pmod h$.
Then there are infinitely many primes $s$ which split
completely in $L$ and satisfy $s\equiv u\pmod M$.
Indeed, the intersection $L\cap\Q(\mu_M)$ is an
abelian subextension of $L/\Q$, and (A.44.4) makes it
exactly $K$. The restriction map is an isomorphism
\[
 \operatorname{Gal}(L\Q(\mu_M)/\Q)
 \xrightarrow{\ \sim\ }
 \operatorname{Gal}(L/\Q)
 \mathop{\times}_{\operatorname{Gal}(K/\Q)}
 \operatorname{Gal}(\Q(\mu_M)/\Q).
 \tag{A.44.5}
\]
The pair $(1,u)$ belongs to this fibre product because
both restrictions to $K$ are the identity. Chebotarev's
theorem \cite[Theorem 8.31]{AgentMilneANT} supplies
infinitely many primes with this Frobenius. Any finite
set of primes can be excluded. This proves the needed
compatibility of the splitting and congruence conditions.

\emph{An elementary bound for the remaining primes.}
Let $p$ be prime, let $p\nmid m a b$, and put
$d=\gcd(m,p-1)$. The number $N_p$ of pairs
$(X,Y)\in(\mathbf F_p^\times)^2$ with
$aX^m+bY^m=1$ satisfies
\[
 N_p\ge p-2-3(d-1)-(d-1)(d-2)\sqrt p.
 \tag{A.44.6}
\]
For completeness, let $\mathcal X_d$ be the $d$
multiplicative characters of $\mathbf F_p^\times$
whose orders divide $d$, all extended by zero at zero,
including the trivial character. Orthogonality gives
\[
 N_p=\sum_{\chi,\psi\in\mathcal X_d}
   \chi(a^{-1})\psi(b^{-1})J(\chi,\psi),
 \qquad
 J(\chi,\psi)=\sum_{t\in\mathbf F_p}
                      \chi(t)\psi(1-t).
\]
The term with both characters trivial is $p-2$.
There are $3(d-1)$ terms of absolute value one:
one character trivial and the other nontrivial, or
both nontrivial with product trivial. Their values
follow from summing a nontrivial character and from
the substitution $t\mapsto t/(1-t)$.
The remaining $(d-1)(d-2)$ terms have absolute value
$\sqrt p$. To see this directly, for nontrivial $\chi$
put $G(\chi)=\sum_{t\ne0}\chi(t)e^{2\pi it/p}$.
Changing variables $t=uv$ in $|G(\chi)|^2$ gives
$|G(\chi)|^2=p$: the inner sum over $v\ne0$ is
$p-1$ for $u=1$ and $-1$ otherwise. When
$\chi\psi$ is nontrivial, changing variables
$(x,y)=(zt,z(1-t))$ gives
$G(\chi)G(\psi)=J(\chi,\psi)G(\chi\psi)$.
These identities prove (A.44.6).

Set $B(m)=(m+1)^4$. If $p>B(m)$, replace $d$ by
$m$ in the lower bound and put $T=\sqrt p$.
The polynomial
$T^2-(m-1)(m-2)T-3m+1$ is increasing for
$T\ge(m+1)^2$, and its value at $(m+1)^2$ is
$m^2(5m+9)>0$. Thus $N_p>0$.
At any such point the derivative in at least one
coordinate is nonzero modulo $p$, so Hensel lifting
gives a point in $(\Z_p^\times)^2$.

We will also use the following precise Hensel estimate.
For $f\in\Z_p[T]$, if
$v_p(f(1))>2v_p(f'(1))=2e$, there is a root $z$
with
\[
                 v_p(z-1)\ge v_p(f(1))-e.
 \tag{A.44.7}
\]
Newton iteration proves the stated bound: its first
increment has valuation $v_p(f(1))-e>e$; the Taylor
remainder has at least twice that valuation; the
derivative keeps valuation $e$ and the successive
increments have increasing valuation. They converge
to the required root and retain the first-increment
bound. The simple-root and Newton forms are also
\cite[Proposition 7.31 and Theorem 7.32]{AgentMilneANT}.

\emph{The induction and its invariant.}
For $j\ge0$ put
\[
 n_j=(j+4)!,\qquad
 H_j(p)=2v_p(n_{j+1})+j+3.
 \tag{A.44.8}
\]
We construct nonzero integers $a_j,b_j$ satisfying
the following three requirements:
\[
 \begin{gathered}
 a_j=3u_j^8,\quad b_j=-2v_j^8,\quad u_j,v_j>0,\\
 \gcd(u_jv_j,6)=1,\qquad \gcd(a_j,b_j)=1;
 \end{gathered}
 \tag{A.44.9}
\]
\[
 a_jX^{n_j}+b_jY^{n_j}=1
 \quad\text{has a point in }(\Z_p^\times)^2
       \text{ for every finite }p;
 \tag{A.44.10}
\]
\[
 v_p(a_j+b_j-1)\ge H_j(p)
       \quad\text{for every }p\le B(n_j).
 \tag{A.44.11}
\]
Initially take $a_0=3$, $b_0=-2$ and
$u_0=v_0=1$. The point $(1,1)$ and the equality
$a_0+b_0=1$ verify all three requirements.

Suppose the construction has reached $j$. Write
\[
 n=n_j,\quad m=n_{j+1},\quad a=a_j,\quad b=b_j,
 \qquad
 S=\{p:p\le B(m)\}\cup\{p:p\mid ab\}.
 \tag{A.44.12}
\]
This is a finite set containing every prime dividing
$m$. For $p\in S$, put
\[
 L_p=H_{j+1}(p)=2v_p(n_{j+2})+j+4.
 \tag{A.44.13}
\]
In particular $L_p>2v_p(m)$ and
$L_p\ge v_p(m)$.

Choose a unit-coordinate solution $(\xi_p,\eta_p)$
of the old equation $aX^n+bY^n=1$ in $\Q_p$.
At primes dividing $m$ we make a specified additional
choice:
\[
        \xi_p\equiv\eta_p\equiv1\pmod {p^{v_p(m)}}.
 \tag{A.44.14}
\]
Here is why this is available. Every prime divisor of
$m=(j+5)!$ is at most $j+5$, hence at most $B(n)$.
Thus (A.44.11) applies. At least one of $a,b$ is a
$p$-adic unit by (A.44.9). Keep the variable of the
other coefficient equal to one and apply (A.44.7) to
the remaining polynomial at one. With
$e=v_p(n)$ and $d=v_p(m)$, its error has valuation
at least $H_j(p)=2d+j+3>2e$. The chosen root differs
from one by an element of valuation at least
$H_j(p)-e\ge d$. This proves (A.44.14), including
$p=2$. At other primes of $S$, (A.44.10) supplies
the required unit solution.

Set
\[
                  M=\prod_{p\in S}p^{L_p}.
\]
This modulus is divisible by $m$. The Chinese remainder
theorem specifies a unit class $Q_0\pmod M$ by
$Q_0\equiv\xi_p\pmod {p^{L_p}}$ for all $p\in S$.
Equation (A.44.14) gives $Q_0\equiv1\pmod m$.
Apply (A.44.5), first with $h=n$, $c=b$, and this
modulus $M$. We obtain a prime $q$ outside $S$, larger
than every prime in $S$ and than $n_{j+2}$, with
\[
 q\equiv\xi_p\pmod {p^{L_p}}\quad(p\in S),
 \qquad q\text{ splits completely in }
                    \Q(\mu_n,b^{1/n}).
 \tag{A.44.15}
\]
At $q$, the number $b$ is an $n$th-power unit. Since
$q\nmid n$, Hensel lifting gives a unit $z_q\in\Z_q$
with
\[
                   b z_q^n=1-aq^n.
 \tag{A.44.16}
\]
We only need its residue modulo $q$ in the next step.

The new positive coefficient $a q^n$ still has the
form $3w^8$ with $\gcd(w,6)=1$, since $8\mid n$
and $q\ne2,3$. The Chinese remainder theorem specifies
a unit class $R_0\pmod {Mq}$ by
\[
 R_0\equiv\eta_p\pmod {p^{L_p}}\quad(p\in S),
 \qquad R_0\equiv z_q\pmod q.
\]
It satisfies $R_0\equiv1\pmod m$. Applying (A.44.5)
with $h=m$, $c=a q^n$, and modulus $Mq$, choose a
prime $r>q$, outside $S\cup\{q\}$, satisfying
\[
 r\equiv R_0\pmod {Mq},
 \qquad r\text{ splits completely in }
               \Q(\mu_m,(a q^n)^{1/m}).
 \tag{A.44.17}
\]
Define
\[
 q_j=q,\quad r_j=r,\qquad
 a_{j+1}=a q^n,\quad b_{j+1}=b r^n.
 \tag{A.44.18}
\]
Their disjoint prime supports and their eighth-power
forms give (A.44.9) at the next stage. In particular
$v_3(a_{j+1})=v_2(b_{j+1})=1$ remains exact.

We verify every finite prime in (A.44.10), with new
exponent $m$. For $p\in S$, both coefficients are
integral and (A.44.15)--(A.44.17) imply
\[
 a q^n+b r^n-1\equiv
       a\xi_p^n+b\eta_p^n-1=0\pmod {p^{L_p}}.
 \tag{A.44.19}
\]
One coefficient is a unit. Apply (A.44.7) at
$(X,Y)=(1,1)$, fixing the other variable, with
derivative valuation $v_p(m)$. The inequality
$L_p>2v_p(m)$ gives a unit-coordinate point of the
new equation. Since every $p\le B(m)$ belongs to
$S$, (A.44.19) also proves (A.44.11) at stage $j+1$.

At $p=q$, equations (A.44.16)--(A.44.17) make the
new equation hold at $(1,1)$ modulo $q$; the
derivative in $Y$ is a unit because $q\nmid mbr$.
It therefore has a unit-coordinate lift. At $p=r$,
set $Y=1$. Complete splitting in (A.44.17) makes
$a q^n$ an $m$th power modulo $r$. The equation
\[
                  X^m=\frac{1-b r^n}{a q^n}
\]
has a nonzero residue root, and $r\nmid m$ gives a
unit root in $\Z_r$. Finally, at a prime outside
$S\cup\{q,r\}$ both new coefficients are units,
and the prime is larger than $B(m)$. The finite-field
bound (A.44.6) and simple-root lifting supply the
required unit-coordinate point. These cases exhaust
the finite primes and complete the induction.

\emph{The completed products and their ordinary local values.}
The choices of $q,r$ exclude all previous prime factors,
so all $q_j,r_j$ are pairwise distinct. For each positive
integer $N$, the factorial $n_j$ is divisible by $N$
for all sufficiently large $j$. Equations (A.44.18)
therefore make both $(a_j)$ and $(b_j)$ Cauchy in
$\widehat{\Q^\times}$: their images are eventually
constant in each simultaneous quotient by $N$th
powers. Their limits are exactly (A.44.1), and the
coordinate calculation gives (A.44.2). In particular,
the valuation vectors have finite support modulo every
$N$, as required in A.10.

Fix a finite prime $p$. Its valuation in each sequence
is eventually constant, and every later ratio
$a_{j+1}/a_j$ or $b_{j+1}/b_j$ is an $n_j$th power
of a $p$-adic unit. For each $k$, the finite group
$\Z_p^\times/(1+p^k\Z_p)$ has finite exponent.
The factorials $n_j$ eventually annihilate this
group. Thus the unit parts of both sequences are
Cauchy in $\Z_p^\times$ and converge to units.
Their limits $x_p,y_p\in\Q_p^\times$ have the finite
nonnegative valuations in (A.44.2), so they lie in
$\Z_p\setminus\{0\}$. Compatibility of the maps
\[
 \Q^\times/(\Q^\times)^N
       \longrightarrow
 \Q_p^\times/(\Q_p^\times)^N
\]
with the two limits gives
$A_p=\kappa(x_p)$ and $B_p=\kappa(y_p)$.
Indeed, the local power classes stabilize both by
the global Cauchy condition and by ordinary local
convergence: the subgroup of $N$th powers of the
local unit group is open, by the same Hensel estimate.

Eventually $p\le B(n_j)$, and (A.44.11) then gives
$v_p(a_j+b_j-1)\ge j+3\longrightarrow\infty$.
Taking the ordinary $p$-adic limits proves
$x_p+y_p=1$. Both are nonzero, so $x_p\in U(\Q_p)$.
This proves $(A,B)\in\mathcal G$ at every finite
prime, including two and three. At the real place
their completed Kummer classes are the signs $+,-$,
represented by the ordinary pair $(3,1-3)$.
Finally (A.44.2)
has infinitely many nonzero valuations in each
coordinate. A rational number has only finitely
many, proving the final assertion.

\emph{A fully specified recursive choice.}
The construction can be read as a terminating algorithm
for every finite stage. At $p\mid m$, use Newton
iteration from one in the unit-coefficient variable.
At other $p\in S$, take the least pair of nonzero
residues solving the old equation modulo $p$, and
Hensel-lift a variable with unit derivative, keeping
the other fixed. These are finite calculations modulo
$p^{L_p}$. Select the least eligible prime $q$ above
the specified bound. Its splitting condition is the
finite test
$b^{(q-1)/n}\equiv1\pmod q$, since
$q\equiv1\pmod n$. Select the least nonzero root
of $bZ^n=1$ modulo $q$, form the second CRT class,
and select the least eligible $r$. Its splitting
test is
$(a q^n)^{(r-1)/m}\equiv1\pmod r$.
Chebotarev, with the exact intersection (A.44.5),
proves termination of both prime searches at every
stage. This specifies the two infinite products
without any unproved simultaneous-primality assertion.
\end{proof}

Together with A.27 and A.38, this gives infinitely
supported ordinary-local Kummer pairs whose full
continuous cup vanishes and which lift to continuous
sections of the actual full profinite class-two
quotient of the tripod. Their abelianizations exclude
rational point sections by (A.44.2). They also exclude
cuspidal sections: at $0$ one coordinate has zero
valuation everywhere, at $1$ the other does, and at
$\infty$ the two valuation vectors agree; none of
these occurs in (A.44.2). No lift to the full
fundamental group or to weight three is asserted.
The finite certificates in
\texttt{code/ghost\_construction\_certificate.py}
check the commutator arithmetic, the finite-field
counting bound, and the precision inequalities.
They do not purport to execute the extremely large
first prime search of the displayed recursion.
